\section{AIME II 2000}
        \begin{enumerate}
        	\item (\textit{AIME II 2000}) The number


 $\frac 2{\log_4{2000^6}} + \frac 3{\log_5{2000^6}}$can be written as $\frac mn$ where $m$ and $n$ are relatively prime positive integers.  Find $m + n$.


 
 \textbf{Answer: }007



\textbf{Solution 1}

$\frac 2{\log_4{2000^6}} + \frac 3{\log_5{2000^6}}$

 $=\frac{\log_4{16}}{\log_4{2000^6}}+\frac{\log_5{125}}{\log_5{2000^6}}$

 $=\frac{\log{16}}{\log{2000^6}}+\frac{\log{125}}{\log{2000^6}}$

 $=\frac{\log{2000}}{\log{2000^6}}$

 $=\frac{\log{2000}}{6\log{2000}}$

 $=\frac{1}{6}$

 Therefore, $m+n=1+6=\boxed{007}$

 





\textbf{Solution 2}

Alternatively, we could've noted that, because $\frac 1{\log_a{b}} = \log_b{a}$

 
 \begin{align*} \frac 2{\log_4{2000^6}} + \frac 3{\log_5{2000^6}} &= 2 \cdot \frac{1}{\log_4{2000^6}} + 3\cdot \frac {1}{\log_5{2000^6} }\\ &=2{\log_{2000^6}{4}} + 3{\log_{2000^6}{5}} \\ &={\log_{2000^6}{4^2}} + {\log_{2000^6}{5^3}}\\ &={\log_{2000^6}{4^2 \cdot 5^3}}\\ &={\log_{2000^6}{2000}}\\ &= {\frac{1}{6}}.\end{align*}


 Therefore our answer is $1 + 6 = \boxed{007}$.

 





\textbf{Solution 3}

We know that $2 = \log_4{16}$ and $3 = \log_5{125}$, and by base of change formula, $\log_a{b} = \frac{\log_c{b}}{\log_c{a}}$. Lastly, notice $\log a + \log b = \log ab$ for all bases.
 \begin{align*} \frac 2{\log_4{2000^6}} + \frac 3{\log_5{2000^6}} = \log_{2000^6}{16} + \log_{2000^6}{125} = \log_{2000^6}{2000} = \frac16 \implies \boxed{007} \end{align*}


 $\bold{Solution}$ $\bold{written}$ $\bold{by}$

 





\textbf{Solution 4}


 \[\frac{2}{\log_4 2000^6} + \frac{3}{\log_5 2000^6}\] 
 \[= \frac{1}{3\log_4 2000} + \frac{1}{2\log_5 2000}\] 
 \[= \frac{1}{3} \log_{2000} 4 + \frac{1}{2} \log_{2000} 5\] 
 \[= \log_{2000} (\sqrt[3]{4} \cdot \sqrt{5}) = x\] 
 \[\implies 2^{4x} \cdot 5^{3x} = 2^{\frac{2}{3}} \cdot 5^{\frac{1}{2}}\] 
 \[\implies 4x + (3\log_2 5)x = \frac{2}{3}+\frac{1}{2} \log_2 5\] 
 \[\implies x = \frac{\frac{2}{3} + \frac{1}{2} \log_2 5}{4 + 3\log_2 5}\] 
 \[\implies 6x = \frac{4 + 3\log_2 5}{4 + 3\log_2 5}\] 
 \[\implies x = \frac{1}{6}\] 
 \[\implies m + n = \boxed{007}\]

 




	\item (\textit{AIME II 2000}) A point whose coordinates are both integers is called a lattice point.  How many lattice points lie on the hyperbola $x^2 - y^2 = 2000^2$?


 
 \textbf{Answer: }098



\textbf{Solution}


 \[(x-y)(x+y)=2000^2=2^8 \cdot 5^6\]

 Note that $(x-y)$ and $(x+y)$ have the same \href{https://artofproblemsolving.com/wiki/index.php?title=Parity}{parities}, so both must be even. We first give a factor of $2$ to both $(x-y)$ and $(x+y)$. We have $2^6 \cdot 5^6$ left. Since there are $7 \cdot 7=49$ factors of $2^6 \cdot 5^6$, and since both $x$ and $y$ can be negative, this gives us $49\cdot2=\boxed{098}$ lattice points.

 





\textbf{Solution 2}

As with solution 1, note that both $x-y$ and $x+y$ must have the same parities, meaning both have to be even. Additionally, we can express both of them in terms of $2^a\cdot3^b$ and $2^c\cdot3^d$. Now, $a+c$ must be equal to 6, and both have to be greater than or equal to 1, so there are by stars and bars 7 ways to do this. Similarly, for $b+d$, we have that both only need to be greater than or equal to 0, so this time there are 7 ways to do so. Since both can be negative, we multiply $7\cdot7\cdot2$ which gives $098$.

 





\textbf{Solution 3}

If we restrict ourselves to the first quadrant, this is equivalent to finding \href{https://artofproblemsolving.com/wiki/index.php?title=Pythagorean\_triple}{Pythagorean triples} for $2000^2 + y^2 = x^2$. We know that every Pythagorean triple corresponds to a pair of integers $m$ and $n$ giving:

 
 \begin{align*} y = m^2 - n^2, && b = 2mn, && x = m^2 + n^2 \end{align*}


 If we let $b=2mn=2000$ then each Pythagorean triple corresponds to a factorization $mn = 1000 = 2^35^3$, of which there are $4\times4=25$.

 But we've been only looking at the first quadrant. If we reflect this quadrant to the others, and eliminate the two duplicate reflections where $y=0$, we end up with $25\times4-2 = 098$ solutions.

 




	\item (\textit{AIME II 2000}) A deck of forty cards consists of four 1's, four 2's,..., and four 10's.  A matching pair (two cards with the same number) is removed from the deck.  Given that these cards are not returned to the deck, let $m/n$ be the probability that two randomly selected cards also form a pair, where $m$ and $n$ are relatively prime positive integers.  Find $m + n.$


 
 \textbf{Answer: }758



\textbf{Solution 1}

There are ${38 \choose 2} = 703$ ways we can draw two cards from the reduced deck. The two cards will form a pair if both are one of the nine numbers that were not removed, which can happen in $9{4 \choose 2} = 54$ ways, or if the two cards are the remaining two cards of the number that was removed, which can happen in $1$ way. Thus, the answer is $\frac{54+1}{703} = \frac{55}{703}$, and $m+n = \boxed{758}$.

 





\textbf{Solution 2}

Instead of counting the cases and doing $\frac{\text{cases wanted}}{\text{total amount of cases}}$ we can use probability directly.

 
For sake of simplicity, WLOG, assume that a pair of ones were removed from the deck of forty cards. We can split it into two cases:

 
Case 1: The pair is ones.

 The probability that a one is chosen is $\frac{2}{38}.$ The probability that a second one is chosen is $\frac{1}{37}$ because one card was removed. Therefore, the probability that the pair is ones is $\frac{2}{38} \cdot \frac{1}{37}.$

 Case 2: The pair is $2-10.$

 The probability that any other number is chosen is $\frac{36}{38}.$ The probability that a number that is equal to this number is chosen (for example, if two was chosen originally then another two being chosen) is $\frac{3}{37}.$ Therefore, the probability that the pair is $2-10$ is $\frac{36}{38} \cdot \frac{3}{37}.$

 Adding these two probabilities gives $\frac{2}{38} \cdot \frac{1}{37} + \frac{36}{38} \cdot \frac{3}{37} = \frac{110}{38 \cdot 37} = \frac{55}{703}$, and $m+n = \boxed{758}.$

 




	\item (\textit{AIME II 2000}) What is the smallest positive integer with six positive odd integer divisors and twelve positive even integer divisors?


 
 \textbf{Answer: }180



\textbf{Solution 1}

We use the fact that the number of divisors of a number $n = p_1^{e_1}p_2^{e_2} \cdots p_k^{e_k}$ is $(e_1 + 1)(e_2 + 1) \cdots (e_k + 1)$. If a number has $18 = 2 \cdot 3 \cdot 3$ factors, then it can have at most $3$ distinct primes in its factorization. 

 Dividing the greatest power of $2$ from $n$, we have an odd integer with six positive divisors, which indicates that it either is ($6 = 2 \cdot 3$) a prime raised to the $5$th power, or two primes, one of which is squared. The smallest example of the former is $3^5 = 243$, while the smallest example of the latter is $3^2 \cdot 5 = 45$. 

 Suppose we now divide all of the odd factors from $n$; then we require a power of $2$ with $\frac{18}{6} = 3$ factors, namely $2^{3-1} = 4$. Thus, our answer is $2^2 \cdot 3^2 \cdot 5 = \boxed{180}$.

 





\textbf{Solution 2}

Somewhat similar to the first solution, we see that the number $n$ has two even factors for every odd factor.  Thus, if $x$ is an odd factor of $n$, then $2x$ and $4x$ must be the two corresponding even factors.  So, the prime factorization of $n$ is $2^2 3^a 5^b 7^c...$ for some set of integers $a, b, c, ...$

 Since there are $18$ factors of $n$, we can write:

 $(2+1)(a+1)(b+1)(c+1)... = 18$

 $(a+1)(b+1)(c+1)... = 6$

 Since $6$ only has factors from the set $1, 2, 3, 6$, either $a=5$ and all other variables are $0$, or $a=3$ and $b=2$, with again all other variables equal to $0$.  This gives the two numbers $2^2 \cdot 3^5$ and $2^2 \cdot 3^2 \cdot 5$.  The latter number is smaller, and is equal to $\boxed {180}$.

 





\textbf{Solution 3}

We see that the least number with 6 odd factors is $3^2*5$. Multiplied by $2^2$ (as each factor of 2 doubles the odd factors, as it can be 2n or $2^2n$. Finally, you get $180$

 -dragoon

 




	\item (\textit{AIME II 2000}) Given eight distinguishable rings, let $n$ be the number of possible five-ring arrangements on the four fingers (not the thumb) of one hand. The order of rings on each finger is significant, but it is not required that each finger have a ring. Find the leftmost three nonzero digits of $n$.


 
 \textbf{Answer: }376



\textbf{Solution}

There are $\binom{8}{5}$ ways to choose the rings, and there are $5!$ distinct arrangements to order the rings [we order them so that the first ring is the bottom-most on the first finger that actually has a ring, and so forth]. The number of ways to distribute the rings among the fingers is equivalent the number of ways we can drop five balls into 4 urns, or similarly dropping five balls into four compartments split by three dividers. The number of ways to arrange those dividers and balls is just $\binom {8}{3}$.

 Multiplying gives the answer: $\binom{8}{5}\binom{8}{3}5! = 376320$, and the three leftmost digits are $\boxed{376}$.

 


 





\textbf{Solution 2}

There are $\binom{8}{5}$to choose the rings.Now, call the rings $A$, $B$, $C$, $D$, and $E$. Any possible ring combination can be generating by scrambling $ABCDE$ and adding three dividers - the things to the left of the first divider will be on the first finger, the things between the first and the second will be on the second finger, and so on. For example, $A|D||CBE$ represents ring $A$ on the first finger, $D$ on the second, none on the third, and $C$, $B$, and $E$ in that order on the fourth. In other words, we simply need to count the number of distinguishable rearrangements of $ABCDE|||$. There are $\frac{8!}{3!}$ of them. 

 Thus the total number of ways is $\binom{8}{5}\frac{8!}{3!} = 376320$ and we submit $\boxed{376}$.

 




	\item (\textit{AIME II 2000}) One base of a trapezoid is $100$ units longer than the other base. The segment that joins the midpoints of the legs divides the trapezoid into two regions whose areas are in the ratio $2: 3$. Let $x$ be the length of the segment joining the legs of the trapezoid that is parallel to the bases and that divides the trapezoid into two regions of equal area. Find the greatest integer that does not exceed $x^2/100$.


 
 \textbf{Answer: }181



\textbf{Solution}

Let the shorter base have length $b$ (so the longer has length $b+100$), and let the height be $h$. The length of the midline of the trapezoid is the average of its bases, which is $\frac{b+b+100}{2} = b+50$. The two regions which the midline divides the trapezoid into are two smaller trapezoids, both with height $h/2$. Then,

 
                    \begin{figure}[H]
                        \centering
                        \includegraphics[width=0.5\textwidth]{Solution/AIME_II_2000_6_Solution_Figure_1.png}
                    \end{figure}
                    
 \[\frac{\frac 12 (h/2) (b + b+50)}{\frac 12 (h/2) (b + 50 + b + 100)} = \frac{2}{3} \Longrightarrow \frac{b + 75}{b + 25} = \frac 32 \Longrightarrow b = 75\]

 Construct the perpendiculars from the vertices of the shorter base to the longer base. This splits the trapezoid into a rectangle and two triangles; it also splits the desired line segment into three partitions with lengths $x_1, 75, x_2$. By similar triangles, we easily find that $\frac{x - 75}{100} = \frac{x_1+x_2}{100} = \frac{h_1}{h}$.

 
                    \begin{figure}[H]
                        \centering
                        \includegraphics[width=0.5\textwidth]{Solution/AIME_II_2000_6_Solution_Figure_2.png}
                    \end{figure}
                    The area of the region including the shorter base must be half of the area of the entire trapezoid, so

 
 \[2 \cdot \frac 12 h_1 (75 + x) = \frac 12 h (75 + 175) \Longrightarrow x = 125 \cdot \frac{h}{h_1} - 75\]

 Substituting our expression for $\frac h{h_1}$ from above, we find that

 
 \[x = \frac{12500}{x-75} - 75 \Longrightarrow x^2 - 75x = 5625 + 12500 - 75x \Longrightarrow x^2 = 18125\]

 Note: When you find the ratio of the heights, you can just plug in $(x-75)$ as $h_1$ and $100$ as $h$ because we only care about the ratio- the actual height will get canceled out in the equation relating the two trapezoids. This simplifies the equation directly from the ratio to 
 \[2\cdot\frac{1}{2}(x+75)(x-75)=\frac{1}{2}(75+175)(100)\] which using difference of squares simplifies to $x^2=18125$.

 The answer is $\left\lfloor\frac{x^2}{100}\right\rfloor = \boxed{181}$.

 




	\item (\textit{AIME II 2000}) 
Given that 

 $\frac 1{2!17!}+\frac 1{3!16!}+\frac 1{4!15!}+\frac 1{5!14!}+\frac 1{6!13!}+\frac 1{7!12!}+\frac 1{8!11!}+\frac 1{9!10!}=\frac N{1!18!}$ find the greatest integer that is less than $\frac N{100}$.


 
 \textbf{Answer: }137



\textbf{Solution}

Multiplying both sides by $19!$ yields: 

 
 \[\frac {19!}{2!17!}+\frac {19!}{3!16!}+\frac {19!}{4!15!}+\frac {19!}{5!14!}+\frac {19!}{6!13!}+\frac {19!}{7!12!}+\frac {19!}{8!11!}+\frac {19!}{9!10!}=\frac {19!N}{1!18!}.\]

 
 \[\binom{19}{2}+\binom{19}{3}+\binom{19}{4}+\binom{19}{5}+\binom{19}{6}+\binom{19}{7}+\binom{19}{8}+\binom{19}{9} = 19N.\] 

 Recall the \href{https://artofproblemsolving.com/wiki/index.php?title=Combinatorial\_identity}{Combinatorial Identity} $2^{19} = \sum_{n=0}^{19} {19 \choose n}$. Since ${19 \choose n} = {19 \choose 19-n}$, it follows that $\sum_{n=0}^{9} {19 \choose n} = \frac{2^{19}}{2} = 2^{18}$.

 Thus, $19N = 2^{18}-\binom{19}{1}-\binom{19}{0}=2^{18}-19-1 = (2^9)^2-20 = (512)^2-20 = 262124$.

 So, $N=\frac{262124}{19}=13796$ and $\left\lfloor \frac{N}{100} \right\rfloor =\boxed{137}$.

 


 





\textbf{Solution 2}

Let $f(x) = (1+x)^{19}.$ Applying the binomial theorem gives us $f(x) = \binom{19}{19} x^{19} + \binom{19}{18} x^{18} + \binom{19}{17} x^{17}+ \cdots + \binom{19}{0}.$ Since $\frac 1{2!17!}+\frac 1{3!16!}+\dots+\frac 1{8!11!}+\frac 1{9!10!} = \frac{\frac{f(1)}{2} - \binom{19}{19} - \binom{19}{18}}{19!},$ $N = \frac{2^{18}-20}{19}.$ After some fairly easy bashing, we get $\boxed{137}$ as the answer.

 





\textbf{Solution 3 (Brute Force)}

Convert each denominator to $19!$ and get the numerators to be $9,51,204,612,1428,2652,3978,4862$ (refer to note). Adding these up we have $13796$ therefore $\boxed{137}$ is the desired answer.

 Note:Notice that each numerator is increased each time by a factor of $\frac{17}{3}, \frac{16}{4}, \frac{15}{5}, \frac{14}{6},$ etc. until $\frac{11}{9}$. If you were taking the test under normal time conditions, it shouldn't be too hard to bash out all of the numbers but it is priority to be careful.

 




	\item (\textit{AIME II 2000}) In trapezoid $ABCD$, leg $\overline{BC}$ is perpendicular to bases $\overline{AB}$ and $\overline{CD}$, and diagonals $\overline{AC}$ and $\overline{BD}$ are perpendicular. Given that $AB=\sqrt{11}$ and $AD=\sqrt{1001}$, find $BC^2$.


 
 \textbf{Answer: }110



\textbf{Solution 1}

Let $x = BC$ be the height of the trapezoid, and let $y = CD$. Since $AC \perp BD$, it follows that $\triangle BAC \sim \triangle CBD$, so $\frac{x}{\sqrt{11}} = \frac{y}{x} \Longrightarrow x^2 = y\sqrt{11}$.

 Let $E$ be the foot of the altitude from $A$ to $\overline{CD}$. Then $AE = x$, and $ADE$ is a \href{https://artofproblemsolving.com/wiki/index.php?title=Right\_triangle}{right triangle}. By the \href{https://artofproblemsolving.com/wiki/index.php?title=Pythagorean\_Theorem}{Pythagorean Theorem}, 

 
 \[x^2 + \left(y-\sqrt{11}\right)^2 = 1001 \Longrightarrow x^4 - 11x^2 - 11^2 \cdot 9 \cdot 10 = 0\]

 The positive solution to this \href{https://artofproblemsolving.com/wiki/index.php?title=Quadratic\_equation}{quadratic equation} is $x^2 = \boxed{110}$.

 
                    \begin{figure}[H]
                        \centering
                        \includegraphics[width=0.5\textwidth]{Solution/AIME_II_2000_8_Solution_Figure_1.png}
                    \end{figure}
                    





\textbf{Solution 2}

Let $BC=x$. Dropping the altitude from $A$ and using the Pythagorean Theorem tells us that $CD=\sqrt{11}+\sqrt{1001-x^2}$. Therefore, we know that vector $\vec{BD}=\langle \sqrt{11}+\sqrt{1001-x^2},-x\rangle$ and vector $\vec{AC}=\langle-\sqrt{11},-x\rangle$. Now we know that these vectors are perpendicular, so their dot product is 0.
 \[\vec{BD}\cdot \vec{AC}=-11-\sqrt{11(1001-x^2)}+x^2=0\]
 \[(x^2-11)^2=11(1001-x^2)\]
 \[x^4-11x^2-11\cdot 990=0.\]As above, we can solve this quadratic to get the positve solution $BC^2=x^2=\boxed{110}$.

 





\textbf{Solution 3}

Let $BC=x$ and $CD=y+\sqrt{11}$. From Pythagoras with $AD$, we obtain $x^2+y^2=1001$. Since $AC$ and $BD$ are perpendicular diagonals of a quadrilateral, then $AB^2+CD^2=BC^2+AD^2$, so we have 
 \[\left(y+\sqrt{11}\right)^2+11=x^2+1001.\] Substituting $x^2=1001-y^2$ and simplifying yields 
 \[y^2+\sqrt{11}y-990=0,\] and the quadratic formula gives $y=9\sqrt{11}$. Then from $x^2+y^2=1001$, we plug in $y$ to find $x^2=\boxed{110}$.

 


 





\textbf{Solution 4}

Let $E$ be the intersection of the diagonals. Since the diagonals are perpendicular, applying the Pythagorean Theorem multiple times we have
 \begin{align*} BC^2&=BE^2+CE^2 \\ &=(AB^2-AE^2)+(CD^2-DE^2) \\ &=CD^2+\sqrt{11}^2-(AE^2+DE^2) \\ &=CD^2+11-AD^2 \\ &=CD^2-990 \end{align*}
Followed by dropping the perpendicular like in solution 1, we obtain system of equation
 \[BC^2=CD^2-990\]
 \[BC^2+CD^2-2\sqrt{11}CD=990\]Rearrange the first equation yields
 \[CD^2-BC^2=990\]Equating it with the second equation we have
 \[CD^2-BC^2=BC^2+CD^2-2\sqrt{11}CD\]Which gives $CD=\frac{BC^2}{\sqrt{11}}$.Substituting into equation 1 obtains the quadratic in terms of $BC^2$
 \[(BC^2)^2-11BC^2-11\cdot990=0\]Solving the quadratic to obtain $BC^2=\boxed{110}$.

 




	\item (\textit{AIME II 2000}) Given that $z$ is a complex number such that $z+\frac 1z=2\cos 3^\circ$, find the least integer that is greater than $z^{2000}+\frac 1{z^{2000}}$.


 
 \textbf{Answer: }000



\textbf{Solution}

Using the quadratic equation on $z^2 - (2 \cos 3 )z + 1 = 0$, we have $z = \frac{2\cos 3 \pm \sqrt{4\cos^2 3 - 4}}{2} = \cos 3 \pm i\sin 3 = \text{cis}\,3^{\circ}$.

 There are other ways we can come to this conclusion. Note that if $z$ is on the \href{https://artofproblemsolving.com/wiki/index.php?title=Unit\_circle}{unit circle} in the complex plane, then $z = e^{i\theta} = \cos \theta + i\sin \theta$  and $\frac 1z= e^{-i\theta} = \cos \theta - i\sin \theta$. We have $z+\frac 1z = 2\cos \theta = 2\cos 3^\circ$ and $\theta = 3^\circ$. Alternatively, we could let $z = a + bi$ and solve to get $z=\cos 3^\circ + i\sin 3^\circ$.

 
Using \href{https://artofproblemsolving.com/wiki/index.php?title=De\_Moivre\%27s\_Theorem}{De Moivre's Theorem} we have $z^{2000} = \cos 6000^\circ + i\sin 6000^\circ$, $6000 = 16(360) + 240$, so $z^{2000} = \cos 240^\circ + i\sin 240^\circ$.

 We want $z^{2000}+\frac 1{z^{2000}} = 2\cos 240^\circ = -1$.

 Finally, the least integer greater than $-1$ is $\boxed{000}$.

 





\textbf{Solution 2}

Let $z=re^{i\theta}$. Notice that we have $2\cos(3^{\circ})=e^{i\frac{\pi}{60}}+e^{-i\frac{\pi}{60}}=re^{i\theta}+\frac{1}{r}e^{-i\theta}.$ 

 $r$ must be $1$ (or if you take the magnitude would not be the same). Therefore, $z=e^{i\frac{\pi}{\theta}}$ and plugging into the desired expression, we get $e^{i\frac{100\pi}{3}}+e^{-i\frac{100\pi}{3}}=2\cos{\frac{100\pi}{3}}=-1$. Therefore, the least integer greater is $\boxed{000}.$

 


 





\textbf{Solution 3 Intuitive}

For this solution, we assume that $z^{2000} + 1/z^{2000}$ and $z^{2048} + 1/z^{2048}$ has the same least integer greater than their solution.

 We have $z + 1/z = 2\cos 3$. Since $\cos 3<1$, $2\cos 3<2$. If we square the equation $z + 1/z = 2\cos 3$, we get $z^2 + 2 + 1/(z^2) = 4\cos^2 3$, or $z^2 + 1/(z^2) = 4\cos^2 3 - 2$. $4\cos^2 3 - 2$ is is less than $2$, since $4\cos^2 3$ is less than $4$. If we square the equation again, we get $z^4 + 1/(z^4) = (4\cos^2 3 - 2)^2 -2$.

 Since $4\cos^2 3 - 2$ is less than 2, $(4\cos^2 3 - 2)^2$ is less than 4, and $(4\cos^2 3 - 2)^2 -2$ is less than 2. However $(4\cos^2 3 - 2)^2 -2$ is also less than $4\cos^2 3 - 2$. we can see that every time we square the equation, the right-hand side gets smaller and into the negatives. Since the smallest integer that is allowed as an answer is 0, the smallest integer greater is $\boxed{000}.$

 





\textbf{Solution 4}

First, let $z = a+bi$ where $a$ and $b$ are real numbers. We now have that 
 \[a+bi+\frac{a-bi}{a^2+b^2} = 2 \cos{3^{\circ}}\] given the conditions of the problem. Equating imaginary coefficients, we have that 
 \[b \left( 1 - \frac{1}{a^2+b^2}\right) = 0\] giving us that either $b=0$ or $|z| = 1$. Let's consider the latter case for now. 

 We now know that $a^2+b^2=1$, so when we equate real coefficients we have that $2a = 2 \cos{3^{\circ}}$, therefore $a = \cos{3^{\circ}}$. So, $b = \sin{3^{\circ}}$ and then we can write $z = \text{cis}(3)^{\circ}$.

 By De Moivre's Theorem, 
 \[z^{2000} + \frac{1}{z^{2000}} = \text{cis} (6000)^{\circ} + \text{cis} (-6000)^{\circ}\]. The imaginary parts cancel, leaving us with $2 \cos{6000^{\circ}}$, which is $240 \pmod{360}$. Therefore, it is $-1$, and our answer is $\boxed{000}$. 

 Now, if $b=0$ then we have that $a+\frac{1}{a} = 2 \cos{3^{\circ}}$. Therefore, $a$ is not violating our conditions set above.

 




	\item (\textit{AIME II 2000}) A circle is inscribed in quadrilateral $ABCD$, tangent to $\overline{AB}$ at $P$ and to $\overline{CD}$ at $Q$. Given that $AP=19$, $PB=26$, $CQ=37$, and $QD=23$, find the square of the radius of the circle.


 
 \textbf{Answer: }647



\textbf{Solution 1}

Call the \href{https://artofproblemsolving.com/wiki/index.php?title=Center}{center} of the circle $O$. By drawing the lines from $O$ tangent to the sides and from $O$ to the vertices of the quadrilateral, four pairs of congruent \href{https://artofproblemsolving.com/wiki/index.php?title=Right\_triangle}{right triangles} are formed.

 Thus, $\angle{AOP}+\angle{POB}+\angle{COQ}+\angle{QOD}=180$, or $(\arctan(\frac{19}{r})+\arctan(\frac{26}{r}))+(\arctan(\frac{37}{r})+\arctan(\frac{23}{r}))=180$.

 Take the $\tan$ of both sides and use the identity for $\tan(A+B)$ to get 
 \[\tan(\arctan(\frac{19}{r})+\arctan(\frac{26}{r}))+\tan(\arctan(\frac{37}{r})+\arctan(\frac{23}{r}))=n\cdot0=0.\]

 Use the identity for $\tan(A+B)$ again to get 
 \[\frac{\frac{45}{r}}{1-19\cdot\frac{26}{r^2}}+\frac{\frac{60}{r}}{1-37\cdot\frac{23}{r^2}}=0.\]

 Solving gives $r^2=\boxed{647}$.

 Note: the equation may seem nasty at first, but once you cancel the $r$s and other factors, you are just left with $r^2$. That gives us $647$ quite easily.

 





\textbf{Solution 2}

Just use the area formula for tangential quadrilaterals. The numbers are really big. A terrible problem to work on ($a, b, c,$ and $d$ are the tangent lengths, not the side lengths).
 \[A = \sqrt{(a+b+c+d)(abc+bcd+cda+dab)} = 105\sqrt{647}\]$r^2=\frac{A^2}{(a+b+c+d)^2} = \boxed{647}$.

 





\textbf{Solution 3 (Smart algebra to make 2 less annoying)}

Using the formulas established in solution 2, one notices:
 \[r^2=\frac{A^2}{(a+b+c+d)^2}\]
 \[r^2=\frac{(a+b+c+d)(abc+bcd+cda+abd) }{(a+b+c+d)^2}\]
 \[r^2=\frac{abc+bcd+acd+abd}{a+b+c+d}\]
 \[r^2=\boxed{647}\]

 which is nowhere near as hard of a calculation. In fact, this is basically the same exact calculation done at the end of solution 1, just with less opportunity to cancel coefficients beforehand.

 





\textbf{Solution 4 (Alternate Solution 2)}

The area of a tangential quadrilateral using side and tangential lengths is $K = \sqrt{(abcd) - (eg-fh)^2}$, where $a$, $b$, $c$, and $d$ are the side lengths, and $e$, $f$, $g$, $h$, the tangential lengths. Then, say $e = 19$, $f= 26$, $g = 37$, and $h = 23$, the lengths $a = 45$, $b = 63$, $c = 60$, and $d = 42$. Note that $K = \sqrt{7144200 - 11025} = \sqrt{7133175} = 105\sqrt{647}$.

 Then, use $K = rs$ with $s = \frac{a+b+c+d}{2} = \frac{210}{2} = 105$. Then, $r$ is $\sqrt{647}$, squared is $\boxed{647}$.

 




	\item (\textit{AIME II 2000}) The coordinates of the vertices of isosceles trapezoid $ABCD$ are all integers, with $A=(20,100)$ and $D=(21,107)$. The trapezoid has no horizontal or vertical sides, and $\overline{AB}$ and $\overline{CD}$ are the only parallel sides. The sum of the absolute values of all possible slopes for $\overline{AB}$ is $m/n$, where $m$ and $n$ are relatively prime positive integers. Find $m+n$.


 
 \textbf{Answer: }131



\textbf{Solution}

For simplicity, we translate the points so that $A$ is on the origin and $D = (1,7)$. Suppose $B$ has integer coordinates; then $\overrightarrow{AB}$ is a \href{https://artofproblemsolving.com/wiki/index.php?title=Vector}{vector} with integer parameters (vector knowledge is not necessary for this solution). We construct the \href{https://artofproblemsolving.com/wiki/index.php?title=Perpendicular}{perpendicular} from $A$ to $\overline{CD}$, and let $D' = (a,b)$ be the reflection of $D$ across that perpendicular. Then $ABCD'$ is a \href{https://artofproblemsolving.com/wiki/index.php?title=Parallelogram}{parallelogram}, and $\overrightarrow{AB} = \overrightarrow{D'C}$. Thus, for $C$ to have integer coordinates, it suffices to let $D'$ have integer coordinates.<sup>[\href{https://artofproblemsolving.com\#note\_1}{1}]</sup> 

 
                    \begin{figure}[H]
                        \centering
                        \includegraphics[width=0.5\textwidth]{Solution/AIME_II_2000_11_Solution_Figure_1.png}
                    \end{figure}
                    Let the slope of the perpendicular be $m$. Then the \href{https://artofproblemsolving.com/wiki/index.php?title=Midpoint}{midpoint} of $\overline{DD'}$ lies on the line $y=mx$, so $\frac{b+7}{2} = m \cdot \frac{a+1}{2}$. Also, $AD = AD'$ implies that $a^2 + b^2 = 1^2 + 7^2 = 50$. Combining these two equations yields

 
 \[a^2 + \left(7 - (a+1)m\right)^2 = 50\]

 Since $a$ is an integer, then $7-(a+1)m$ must be an integer. There are $12$ pairs of integers whose squares sum up to $50,$ namely $( \pm 1, \pm 7), (\pm 7, \pm 1), (\pm 5, \pm 5)$. We exclude the cases $(\pm 1, \pm 7)$ because they lead to degenerate parallelograms (rectangle, line segment, vertical and horizontal sides). Thus we have

 
 \[7 - 8m = \pm 1, \quad 7 + 6m = \pm 1, \quad 7 - 6m = \pm 5, 7 + 4m = \pm 5\]

 These yield $m = 1, \frac 34, -1, -\frac 43, 2, \frac 13, -3, - \frac 12$. Therefore, the corresponding slopes of $\overline{AB}$ are $-1, -\frac 43, 1, \frac 34, -\frac 12, -3, \frac 13$, and $2$.  The sum of their absolute values is $\frac{119}{12}$. The answer is $m+n= \boxed{131}$

 
\href{https://artofproblemsolving.com\#ref\_1}{\textasciicircum{}} <small>In other words, since $ABCD'$ is a parallelogram, the difference between the x-coordinates and the y-coordinates of $C$ and $D'$ are, respectively, the difference between the x-coordinates and the y-coordinates of $A$ and $B$. But since the latter are integers, then the former are integers also, so $C$ has integer coordinates \href{https://artofproblemsolving.com/wiki/index.php?title=Iff}{iff} $D'$ has integer coordinates.</small>

 
Note: It may not seem like we considered that the coordinates of $B$ have to be integers, but since the slopes of $\overline{AB}$ are all rational, we can just extend $\overline{AB}$ by some arbitrary amount so that $B$ becomes a lattice point, which won't affect the position of $D'$. ~inaccessibles

 





\textbf{Solution 2}

A very natural solution: Shift $A$ to the origin. Suppose point $B$ is $(x, kx)$. Note $k$ is the slope we're looking for. Note that point $C$ must be of the form:$(x \pm 1, kx \pm 7)$ or $(x \pm 7, kx \pm 1)$ or $(x \pm 5, kx \pm 5)$. Note that we want the slope of the line connecting $D$ and $C$ so also be $k$, since $AB$ and $CD$ are parallel.Instead of dealing with the 12 cases, we consider point $C$ of the form $(x \pm Y, kx \pm Z)$ where we plug in the necessary values for $Y$ and $Z$ after simplifying.Since the slopes of $AB$ and $CD$ must both be $k$,$\frac{7 - kx \pm Z}{1 - x \pm Y} = k \implies k = \frac{7 \pm Z}{1 \pm Y}$. Plugging in the possible values of $\pm 7, \pm 1, \pm 5$ in their respective pairs and ruling out degenerate cases, we find the sum is $\frac{119}{12} \implies m + n = \boxed{131}$- whatRthose

 [Note: This Solution is a lot faster if you rule out $(Y, Z) = (1, 7)$ due to degeneracy: $k$ cannot equal $0$ (parallel to x-axis), undefined (parallel to y-axis), or $7$ (parallel to $AD$).]

 




	\item (\textit{AIME II 2000}) The points $A$, $B$ and $C$ lie on the surface of a sphere with center $O$ and radius $20$. It is given that $AB=13$, $BC=14$, $CA=15$, and that the distance from $O$ to triangle $ABC$ is $\frac{m\sqrt{n}}k$, where $m$, $n$, and $k$ are positive integers, $m$ and $k$ are relatively prime, and $n$ is not divisible by the square of any prime. Find $m+n+k$.


 
 \textbf{Answer: }118



\textbf{Solution 1}

Let $D$ be the foot of the \href{https://artofproblemsolving.com/wiki/index.php?title=Perpendicular}{perpendicular} from $O$ to the plane of $ABC$. By the \href{https://artofproblemsolving.com/wiki/index.php?title=Pythagorean\_Theorem}{Pythagorean Theorem} on triangles $\triangle OAD$, $\triangle OBD$ and $\triangle OCD$ we get:

 
 \[DA^2=DB^2=DC^2=20^2-OD^2\]

 It follows that $DA=DB=DC$, so $D$ is the \href{https://artofproblemsolving.com/wiki/index.php?title=Circumcenter}{circumcenter} of $\triangle ABC$.

 By \href{https://artofproblemsolving.com/wiki/index.php?title=Heron\%27s\_Formula}{Heron's Formula} the area of $\triangle ABC$ is (alternatively, a $13-14-15$ triangle may be split into $9-12-15$ and $5-12-13$ \href{https://artofproblemsolving.com/wiki/index.php?title=Right\_triangle}{right triangles}):

 
 \[K = \sqrt{s(s-a)(s-b)(s-c)} = \sqrt{21(21-15)(21-14)(21-13)} = 84\]

 From $R = \frac{abc}{4K}$, we know that the \href{https://artofproblemsolving.com/wiki/index.php?title=Circumradius}{circumradius} of $\triangle ABC$ is:

 
 \[R = \frac{abc}{4K} = \frac{(13)(14)(15)}{4(84)} = \frac{65}{8}\]

 Thus by the \href{https://artofproblemsolving.com/wiki/index.php?title=Pythagorean\_Theorem}{Pythagorean Theorem} again,

 
 \[OD = \sqrt{20^2-R^2} = \sqrt{20^2-\frac{65^2}{8^2}} = \frac{15\sqrt{95}}{8}.\]

 So the final answer is $15+95+8=\boxed{118}$.

 





\textbf{Solution 2 (Vectors)}

We know the radii to $A$,$B$, and $C$ form a triangular pyramid $OABC$. We know the lengths of the edges $OA = OB = OC = 20$. First we can break up $ABC$ into its two component right triangles $5-12-13$ and $9-12-15$. Let the $y$ axis be perpendicular to the base and $x$ axis run along $BC$, and $z$ occupy the other dimension, with the origin as $C$. We look at vectors $OA$ and $OC$. Since $OAC$ is isoceles we know the vertex is equidistant from $A$ and $C$, hence it is $7$ units along the $x$ axis. Hence for vector $OC$, in form $<x,y,z>$ it is $<7, h, l>$ where $h$ is the height (answer) and $l$ is the component of the vertex along the $z$ axis. Now on vector $OA$, since $A$ is $9$ along $x$, and it is $12$ along $z$ axis, it is $<-2, h, 12- l>$. We know both vector magnitudes are $20$. Solving for $h$ yields $\frac{15\sqrt{95} }{8}$, so Answer = $\boxed{118}$.

 Note: $OA$ is actually $\left< 2, -h, 12-l \right>$. ~fermat\_sLastAMC

 





\textbf{Solution 3}

The distance from $O$ to $ABC$ forms a  right triangle with the circumradius of the triangle and the radius of the sphere. 

 
The hypotenuse has length $20$, since it is the radius of the sphere.

 
The circumradius of a $13$, $14$, $15$ triangle is $\frac{65}{8}$, so the distance from $O$ to $ABC$ is given by $\sqrt{20^2-(\frac{65}{8})^2} = \frac{15\sqrt{95}}{8}$, and $15+95+8 = \boxed{118}$.

 -skibbysiggy

 




	\item (\textit{AIME II 2000}) The equation $2000x^6+100x^5+10x^3+x-2=0$ has exactly two real roots, one of which is $\frac{m+\sqrt{n}}r$, where $m$, $n$ and $r$ are integers, $m$ and $r$ are relatively prime, and $r>0$. Find $m+n+r$.


 
 \textbf{Answer: }200



\textbf{Solution}

We may factor the equation as:<sup>[\href{https://artofproblemsolving.com\#note\_1}{1}]</sup>

 
 \begin{align*} 2000x^6+100x^5+10x^3+x-2&=0\\ 2(1000x^6-1) + x(100x^4+10x^2+1)&=0\\ 2[(10x^2)^3-1]+x[(10x^2)^2+(10x^2)+1]&=0\\ 2(10x^2-1)[(10x^2)^2+(10x^2)+1]+x[(10x^2)^2+(10x^2)+1]&=0\\ (20x^2+x-2)(100x^4+10x^2+1)&=0\\ \end{align*}


 Now $100x^4+10x^2+1\ge 1>0$ for real $x$. Thus the real roots must be the roots of the equation $20x^2+x-2=0$. By the \href{https://artofproblemsolving.com/wiki/index.php?title=Quadratic\_formula}{quadratic formula} the roots of this are:

 
 \[x=\frac{-1\pm\sqrt{1^2-4(-2)(20)}}{40} = \frac{-1\pm\sqrt{1+160}}{40} = \frac{-1\pm\sqrt{161}}{40}.\]

 Thus the only root of the required form is $x=\frac{-1+\sqrt{161}}{40}$, and so the final answer is $-1+161+40 = \boxed{200}$.

 Another Note: One could have also determined that the imaginary roots lie with the polynomial $100x^4+10x^2+1$ through Descartes Rule of Signs.

 
\href{https://artofproblemsolving.com\#ref\_1}{\textasciicircum{}} A well-known technique for dealing with symmetric (or in this case, nearly symmetric) polynomials is to divide through by a power of $x$ with half of the polynomial's degree (in this case, divide through by $x^3$), and then to use one of the substitutions $t = x \pm \frac{1}{x}$. In this case, the substitution $t = x\sqrt{10} - \frac{1}{x\sqrt{10}}$ gives $t^2 + 2 = 10x^2 + \frac 1{10x^2}$ and $2\sqrt{10}(t^3 + 3t) = 200x^3 - \frac{2}{10x^3}$, which reduces the polynomial to just $(t^2 + 3)\left(2\sqrt{10}t + 1\right) = 0$. Then one can backwards solve for $x$.Note: After dividing the equation with $x^3$, divide again with $10$ before substituting it with $t$ to get it right.

 A slightly different approach using symmetry:
Let $y = 10x - 1/x$.Notice that the equation can be rewritten (after dividing across by $x^3$) as 
 \[2( (10x)^3 - 1/x^3 ) + (10x)^2 + (1/x)^2 + 10 = 0\]Now it is easy to see that the equation reduces to 
 \begin{align*} 2(y^3+30y)+ (y^2+20) + 10 = 0 \\ 2y^3 + y^2 + 60y + 30 = 0 \\ y^2(2y+1) + 30(2y+1) = 0 \\ (2y+1)(y^2+30)= 0 \\ \end{align*}
so for real solutions we have $y = -1/2$. Solve the quadratic in $x$ to get the final answer as $\boxed{200}$.

 





\textbf{Solution 2 (Complex Bash)}

It would be really nice if the coefficients were symmetrical. What if we make the substitution $x = -\frac{i}{\sqrt{10}}y$? Then the equation becomes 

 $-2y^6 - \left(\frac{i}{\sqrt{10}}\right)y^5 + \left(\frac{i}{\sqrt{10}}\right)y^3 - \left(\frac{i}{\sqrt{10}}\right)y - 2 = 0$.

 It's symmetric! Dividing by $y^3$ and rearranging, we get 

 $-2\left(y^3 + \frac{1}{y^3}\right) - \left(\frac{i}{\sqrt{10}}\right)\left(y^2 + \frac{1}{y^2}\right) + \left(\frac{i}{\sqrt{10}}\right)$

 Now, if we let $z = y + \frac{1}{y}$, we can get the equations 

 $z = y + \frac{1}{y}$ 

 $z^2 - 2 = y^2 + \frac{1}{y^2}$

 and

 $z^3 - 3z = y^3 + \frac{1}{y^3}$

 (These come from squaring $z$ and subtracting $2$ then multiplying that result by $z$ and subtracting $z$.) Plugging this into our polynomial, expanding, and rearranging, we get 

 $-2z^3 - \left(\frac{i}{\sqrt{10}}\right)z^2 + 6z + \left(\frac{3i}{\sqrt{10}}\right)$

 Now, we see that the two $i$ terms must cancel in order to get this polynomial equal to $0$, so what squared equals 3? Plugging in $z = \sqrt{3}$ into the polynomial, we see that it works! Is there something else that equals 3 when squared? Trying $z = -\sqrt{3}$, we see that it also works! Great, we use long division on the polynomial by 

 $\left(z - \sqrt{3}\right)\left(z + \sqrt{3}\right) = z^2 - 3$ and we get 

 $2z - \left(\frac{i}{\sqrt{10}}\right) = 0$. 

 We know that the other two solutions for z wouldn't result in real solutions for $x$ since we have to solve a quadratic with a negative discriminant and then multiply by $-\left(\frac{i}{\sqrt{10}}\right)$. We get that $z = \left(\frac{i}{-2\sqrt{10}}\right)$. Solving for $y$ (using $y + \frac{1}{y} = z$) we get that $y = \frac{-i \pm \sqrt{161}i}{4\sqrt{10}}$, and multiplying this by $-\left(\frac{i}{\sqrt{10}}\right)$ (because $x = -\left(\frac{i}{\sqrt{10}}\right)y$) we get that $x = \frac{-1 \pm \sqrt{161}}{40}$ for a final answer of $-1 + 161 + 40 = \boxed{200}.$

 





\textbf{Solution 3 (Geometric Series)}

Observe that the given equation may be rearranged as$2000x^6-2+(100x^5+10x^3+x)=0$.The expression in parentheses is a geometric series with common factor $10x^2$. Using the geometric sum formula, we rewrite as$2000x^6-2+\frac{1000x^7-x}{10x^2-1}=0, 10x^2-1\neq0$.Factoring a bit, we get$2(1000x^6-1)+(1000x^6-1)\frac{x}{10x^2-1}=0, 10x^2-1\neq0 \implies$$(1000x^6-1)(2+\frac{x}{10x^2-1})=0, 10x^2-1\neq0$.Note that setting $1000x^6-1=0$ gives $10x^2-1=0$, which is clearly extraneous.Hence, we set $2+\frac{x}{10x^2-1}=0$ and use the quadratic formula to get the desired root$x=\frac{-1+\sqrt{161}}{40} \implies -1+161+40=\boxed{200}$

 





\textbf{Solution 4}

If we look at the polynomial's terms, we can see that the number of zeros in a term more or less correlates to the power of $x^2$. Thus, we let $y=10x^2$. The equation then becomes $2y^3+xy^2+xy+x-2=0$, or $x(y^2+y+1)=2(1-y^3)$.

 By the difference of cubes formula, $2(1-y^3)=2(1-y)(1+y+y^2)$, so we have two cases: either $y^2+y+1=0$, or $x=2(1-y)$. We start with the second formula as it is simpler.

 Solving with the quadratic formula after re-substitution, we see that $x=\frac{-1\pm\sqrt{161}}{40}$, so the answer is $-1+161+40=\boxed{200}$.

 For the sake of completeness, if we check the other equation, we come to the conclusion that $y=\frac{-1\pm i\sqrt{3}}{2}$, so no real solution exists for $x$. Thus our solution is correct.

 





\textbf{Solution 5 (substitutions)}

Let $x = \frac{y}{10}.$ We have 

 
 \[\frac{2000y^6}{10^6} + \frac{100y^5}{10^5}+\frac{10y^3}{10^3}+\frac{y}{10} - 2=0\]
 \[= \frac{2y^6}{10^3}+\frac{y^5}{10^3}+\frac{y^3}{10^2}+\frac{y}{10}-2\]
 \[\Longleftrightarrow 2y^6 + y^5 + 10y^3 + 100y - 2000 = 0.\]

 Divide this by $y^3,$ and we get

 
 \[2y^3 + y^2 + 10 + \frac{100}{y^2} - \frac{2000}{y^3}.\]

 By inspection, this looks like it could be expressed as a polynomial in $z=y - \frac{10}{y}.$First, let's express $2\left(y^3 - \frac{1000}{y^3}\right)$ in terms of $z.$ We have 

 
 \[z^3 = \left(y-\frac{10}{y}\right)^3 = y^3 - 30y + \frac{300}{y} - \frac{1000}{y^3},\]
 \[\Longleftrightarrow z^3 + 30z = \left(y-\frac{10}{y}\right)^3 + 30\left(y - \frac{10}{y} \right) = y^3 - \frac{1000}{y^3}.\]

 So,

 
 \[2y^3 - \frac{2000}{y^3} = 2z^3 + 60z.\]

 Now we want to express $y^2 + 10 + \frac{100}{y^2}$ in terms of $z$ too.

 
 \[z^2 + 30 = y^2 - 20 + \frac{100}{y^2}+30 = y^2 + 10 + \frac{100}{y^2}.\]

 Finally,

 
 \[2y^3 + y^2 + 10 + \frac{100}{y^2} - \frac{2000}{y^3} = 2z^3 + z^2 + 60z + 30.\]

 It is quite easy to see that $-\frac{1}{2}$ is one of the roots of this polynomial. Since the problem specifies that there are exactly two real roots, we know that the other roots of this cubic are nonreal. Now, we backtrack.
 \[z = -\frac{1}{2} = y-\frac{10}{y}.\]
 \[y^2 + \frac{y}{2} - 10 = 0\]
 \[2y^2 + y - 20 = 0\]
 \[\frac{-1 \pm \sqrt{1+160 }}{4}\]
 \[\Longleftrightarrow y = \frac{-1 + \sqrt{161}}{4}.\]Since $x=\frac{y}{10}, x = \frac{-1 + \sqrt{161}}{40},$ which means our answer is 
 \[-1 + 161 + 40 = \boxed{200}.\]

 




	\item (\textit{AIME II 2000}) Every positive integer $k$ has a unique factorial base expansion $(f_1,f_2,f_3,\ldots,f_m)$, meaning that $k=1!\cdot f_1+2!\cdot f_2+3!\cdot f_3+\cdots+m!\cdot f_m$, where each $f_i$ is an integer, $0\le f_i\le i$, and $0<f_m$. Given that $(f_1,f_2,f_3,\ldots,f_j)$ is the factorial base expansion of $16!-32!+48!-64!+\cdots+1968!-1984!+2000!$, find the value of $f_1-f_2+f_3-f_4+\cdots+(-1)^{j+1}f_j$.


 
 \textbf{Answer: }495



\textbf{Solution 1}

Note that $1+\sum_{k=1}^{n-1} {k\cdot k!} = 1+\sum_{k=1}^{n-1} {((k+1) \cdot k!- k!)} = 1+\sum_{k=1}^{n-1} {((k+1)!- k!)}$

 $= 1 + ((2! - 1!) + (3! - 2!) + \cdots + (n! - (n-1)!)) = n!$.

 Thus for all $m\in\mathbb{N}$,

 $(32m+16)!-(32m)! = \left(1+\sum_{k=1}^{32m+15} {k\cdot k!}\right)-\left(1+\sum_{k=1}^{32m-1} {k\cdot k!}\right) = \sum_{k=32m}^{32m+15}k\cdot k!.$

 So now,

 
 \begin{align*} 16!-32!+48!-64!+\cdots+1968!-1984!+2000!&=16!+(48!-32!)+(80!-64!)\cdots+(2000!-1984!)\\ &=16! +\sum_{m=1}^{62}\left((32m+16)!-(32m)!\right)\\ &=16! +\sum_{m=1}^{62}\sum_{k=32m}^{32m+15}k\cdot k! \end{align*}


 Therefore we have $f_{16} = 1$, $f_k=k$ if $32m\le k \le 32m+15$ for some $m=1,2,\ldots,62$, and $f_k = 0$ for all other $k$.

 Therefore we have:

 
 \begin{align*} f_1-f_2+f_3-f_4+\cdots+(-1)^{j+1}f_j &= (-1)^{17}\cdot 1 + \sum_{m=1}^{62}\sum_{k=32m}^{32m+15}(-1)^{k+1}k\\ &= -1 + \sum_{m=1}^{62}\left[\sum_{j=16m}^{16m+7}(-1)^{2j+1}2j+\sum_{j=16m}^{16m+7}(-1)^{2j+2}(2j+1)\right]\\ &= -1 + \sum_{m=1}^{62}\sum_{j=16m}^{16m+7}[(-1)^{2j+1}2j+(-1)^{2j+2}(2j+1)]\\ &= -1 + \sum_{m=1}^{62}\sum_{j=16m}^{16m+7}[-2j+(2j+1)]\\ &= -1 + \sum_{m=1}^{62}\sum_{j=16m}^{16m+7}1\\ &= -1 + \sum_{m=1}^{62}8\\ &= -1 + 8\cdot 62\\ &= \boxed{495}. \end{align*}


 





\textbf{Solution 2 (informal)}

This is equivalent to Solution 1. I put up this solution merely for learners to see the intuition.

 Let us consider a base $n$ number system. It’s a well known fact that when we take the difference of two integral powers of $n$, (such as $10000_{10} - 100_{10}$) the result will be an integer in base $n$ composed only of the digits $n - 1$ and $0$ (in this example, $9900$). More specifically, the difference $(n^k)_n - (n^j)_n$, $j<k$ , is an integer $k$ digits long (note that $(n^k)_n$ has $k + 1$ digits). This integer is made up of $(k-j)$ $(n - 1)$’s followed by $j$ $0$’s.

 It should make sense that this fact carries over to the factorial base, albeit with a modification. Whereas in the general base $n$, the largest digit value is $n - 1$, in the factorial base, the largest digit value is the argument of the factorial in that place. (for example, $321_!$ is a valid factorial base number, as is $3210_!$. However, $31_!$ is not, as $3$ is greater than the argument of the second place factorial, $2$. $31_!$ should be represented as $101_!$, and is $7_{10}$.) Therefore, for example, $1000000_! - 10000_!$ is not $990000_!$, but rather is $650000_!$. Thus, we may add or subtract factorials quite easily by converting each factorial to its factorial base expression, with a $1$ in the argument of the factorial’s place and $0$’s everywhere else, and then using a standard carry/borrow system accounting for the place value. 

 With general intuition about the factorial base system out of the way, we may tackle the problem. We use the associative property of addition to regroup the terms as follows: $(2000! - 1984!) + (1968! - 1952!) + \cdots + (48! - 32!) + 16!$ we now apply our intuition from paragraph 2. $2000!_{10}$ is equivalent to $1$ followed by $1999$ $0$’s in the factorial base, and $1984!$ is $1$ followed by $1983$ $0$’s, and so on. Therefore, $2000! - 1984! = (1999)(1998)(1997)\cdots(1984)$ followed by $1983$ $0$’s in the factorial base. $1968! - 1952! = (1967)(1966)\cdots(1952)$ followed by $1951$ $0$’s, and so on for the rest of the terms, except $16!$, which will merely have a $1$ in the $16!$ place followed by $0$’s. To add these numbers, no carrying will be necessary, because there is only one non-zero value for each place value in the sum. Therefore, the factorial base place value $f_k$ is $k$ for all $32m \leq k \leq 32m+15$ if $1\leq m \in\mathbb{Z} \leq 62$, $f_{16} = 1$, and $f_k = 0$ for all other $k$. 

 Therefore, to answer, we notice that $1999 - 1998 = 1997 - 1996 = 1$, and this will continue. Therefore, $f_{1999} - f_{1998} + \cdots - f_{1984} = 8$. We have 62 sets that sum like this, and each contains $8$ pairs of elements that sum to $1$, so our answer is almost $8 \cdot 62$. However, we must subtract the $1$ in the $f_{16}$ place, and our answer is $8 \cdot 62 - 1 = \boxed{495}$.

 





\textbf{Solution 3 (less formality)}

Let $S = 16!-32!+\cdots-1984!+2000!$. Note that since $|S - 2000!| << 2000!$ (or $|S - 2000!| = 1984! + \cdots$ is significantly smaller than $2000!$), it follows that $1999! < S < 2000!$. Hence $f_{2000} = 0$. Then $2000! = 2000 \cdot 1999! = 1999 \cdot 1999! + 1999!$, and as $S - 2000! << 1999!$, it follows that $1999 \cdot 1999! < S < 2000 \cdot 1999!$. Hence $f_{1999} = 1999$, and we now need to find the factorial base expansion of 

 
 \[S_2 = S - 1999 \cdot 1999! = 1999! - 1984! + 1962! - 1946! + \cdots + 16!\]

 Since $|S_2 - 1999!| << 1999!$, we can repeat the above argument recursively to yield $f_{1998} = 1998$, and so forth down to $f_{1985} = 1985$. Now $S_{16} = 1985! - 1984! + 1962! + \cdots = 1984 \cdot 1984! + 1962! + \cdots$, so $f_{1984} = 1984$. 

 The remaining sum is now just $1962! - 1946! + \cdots + 16!$. We can repeatedly apply the argument from the previous two paragraphs to find that $f_{16} = 1$, and $f_k=k$ if $32m\le k \le 32m+15$ for some $m=1,2,\ldots,62$, and $f_k = 0$ for all other $k$.

 Now for each $m$, we have $-f_{32m} + f_{32m+1} - f_{32m+2} + \cdots + f_{32m + 31}$ $= -32m + (32m + 1) - (32m + 2) + \cdots - (32m - 30) + (32 m + 31)$ $= 1 + 1 + \cdots + 1 + 1$ $= 8$. Thus, our answer is $-f_{16} + 8 \cdot 62 = \boxed{495}$.

 





\textbf{Solution 4 (my brain can't comprehend the other solutions)}

First consider how we would express $48!-32!$. We can't use any multiples of $48!$, since we'll never be able to make the $-32!$. Instead, we have to start with using $47!$s. Having $47\cdot47!$ is the closest we can get to $48!$, since $48! = 48\cdot 47!$. Thus, $f_{47}=47$.

 
Now consider what we have left to express. We wanted to express $48!-32!$, and now we have $47\cdot 47!$. This gives us $47!-32!$ left. Clearly, we can't use anymore $47!$s, so we have to use $46!$s. The most $46!$s we could use is $46$ of them. Thus, $f_{46}=46$. Since $47!=47\cdot 46!$, we are left with $46!-32!$ to express. Clearly, each time, we have $f_n = n$, leaving us with $n!-32!$ to express.

 
Eventually, we will get down to needing to express $34!-32!$. We will need $33$ $33!$s, leaving $33!-32!$ to be expressed. Since $33!=33\cdot 32!$, $33!-32! = 32\cdot 32!$. Thus, $48!-32!$ can be expressed as $\sum_{i=32}^{47} i\cdot (i!)$.

 
Returning to the problem, notice we can break the expression into groups of $2$, having $16!+(48!-32!)+(80!-64!)+...(2000!-1984!)$. Each pair has an alternating sum of $8$, and since there are $62$ pairs, the sum is $496$. Including the $16!$ at the start, which accounts for $-1$, we get $\boxed{495}$.

 -skibbysiggy

 





\textbf{Solution 5 (Much simpler than the other solutions [besides 4])}

Let's solve $3! - 5!$ without numerically solving each. What do we know about number theory? Well, one way to rewrite a subtraction $a$ from $b$ given $a > b$ is to borrow from $a$, add it to $b$, and subtract $a$ as a result. More formally, if $b - a : a > b \rightarrow (b + a) - 2a$.

 How does this apply to factorials? Well, say we borrow from the next highest factorial, that is, if we have $b! - a!$ where $a > b$, then we can borrow once from $(b+1)!$ to obtain $b+1$ copies of $b!$ (which is exactly $(b+1)!$ if you are wondering), and we subtract $(b+1)!$ as a result. So, for our case $3! - 5!$, we can rewrite as $3! - (5)(4)3!$, or $3! - 20 \cdot 3!$, or $-19 \cdot 3!$. We have the constraint however that $K = f_1(1!) + f_2(2!) + ... + f_n(n!)$ where $0 \le f_n \le n$. However, our case swaps the sign of $K$, giving us $0 \ge -f_n \ge -n$. So, we borrow once from $4!$, giving us 4 copies of $3!$, and ultimately resulting in $(-19 + 4)3! - 1 \cdot 4!$. The result $-19 + 4 = -15$ is still negative. To compensate, we take 1 copy of $5!$ to donate to $4!$, which then we can donate the new copies of $4!$ to $3!$. So, giving one copy of $5!$ to $4!$ gives us 5 copies of $4!$, so we have $(-15)3! + (-1 + 5)4! = (-15)3! + 4(4!)$, in which $4 \cdot 4!$ has 16 copies of $3!$ (as $4 \cdot 3! = 4!$). Then, we have $(-15 + 16)3! + 4 \cdot 4! - 1 \cdot 5!$, with a final answer of $3! + 4 \cdot 4! - 1 \cdot 5!$, where $(f_3, f_4, f_5) = (1, 4, -1)$ before normalization (note that -1 is allowed for $f_5$ as the expression itself is negative. For all positive $K$, the inequality must hold). This gives us an idea of what to do for the problem.

 Consider $n! - (n+16)!$. We must borrow 1 copy from $(n+1)!$, which then results in borrowing 2 copies from $(n+2)!$, which then results in borrowing 3 copies from $(n+3)!$, and so on so forth until we borrow 15 copies from $(n+15)!$, and 16 copies from $(n+16)!$. Additionally, the number of copies we borrow from each factorial correspond to the $f_n$ power (We already proved that $f_n$ is determined by the explicit values of the remaining factorial expansion. However, as homework, try to see how each $f_n$ corresponds to the number of copies borrowed). Thus, we have $f_1 - f_2 + f_3, ..., + f_{15} - f_{16} = 1 - 2 + 3 - 4 + ... + 15 - 16 = -8$. This means that each $n! - (n+16)!$ pair contributes -8 to the count. So, it becomes clear that every $(n+16)! - n!$ pair contributes 8 to the count. Thus, we can pair each term as $16! + (48! - 32!) + (...) + (2000! - 1984!)$. There are 62 of the $(n+16)! - n!$ pairs, and thus we see that a total alternating sum of 496 arises. However, we have one last case, that is, $16!$. We borrow 0 copies of $n!$, because 1 copy of $n!$ is already present! Thus, the remaining $(-1)^n (f_n)$ is nothing but $(-1)^{125}(1) = -1$ (there are a total of 125 terms, and 16! is the 125th), and so our final sum is $496 - 1 = \boxed{495}$.

 




	\item (\textit{AIME II 2000}) 
Find the least positive integer $n$ such that 

 $\frac 1{\sin 45^\circ\sin 46^\circ}+\frac 1{\sin 47^\circ\sin 48^\circ}+\cdots+\frac 1{\sin 133^\circ\sin 134^\circ}=\frac 1{\sin n^\circ}.$
 \textbf{Answer: }001



\textbf{Solution 1}

We apply the identity 

 
 \begin{align*} \frac{1}{\sin n \sin (n+1)} &= \frac{1}{\sin 1} \cdot \frac{\sin (n+1) \cos n - \sin n \cos (n+1)}{\sin n \sin (n+1)} \\ &= \frac{1}{\sin 1} \cdot \left(\frac{\cos n}{\sin n} - \frac{\cos (n+1)}{\sin (n+1)}\right) \\ &= \frac{1}{\sin 1} \cdot \left(\cot n - \cot (n+1)\right). \end{align*}


 The motivation for this identity arises from the need to decompose those fractions, possibly into \href{https://artofproblemsolving.com/wiki/index.php?title=Telescoping}{telescoping}. 

 Thus our summation becomes 

 
 \[\sum_{k=23}^{67} \frac{1}{\sin (2k-1) \sin 2k} = \frac{1}{\sin 1} \left(\cot 45 - \cot 46 + \cot 47 - \cdots + \cot 133 - \cot 134 \right).\]

 Since $\cot (180 - x) = - \cot x$, the summation simply reduces to $\frac{1}{\sin 1} \cdot \left( \cot 45 - \cot 90 \right) = \frac{1 - 0}{\sin 1} = \frac{1}{\sin 1^{\circ}}$. Therefore, the answer is $\boxed{001}$.

 





\textbf{Solution 2}

We can make an approximation by observing the following points:

 The average term is around the 60's which gives $\frac{4}{3}$.

 There are 45 terms, so the approximate sum is 60.

 Therefore, $\sin(n^\circ)$ equals approximately $\frac{1}{60}$.

 Recall that the approximation of $\sin(x)$ in radians is x if x is close to zero. In this case x is close to zero. Converting to radians we see that $\sin(1)$ in degrees is about sin$\frac{1}{57}$ in radians, or is about $\frac{1}{57}$ because of the approximation. What we want is apparently close to that so we make the guess that n is equal to 1 degree. Basically, it boils down to the approximation of $\sin(1)=\frac{1}{60}$ in degrees, convert to radians and use the small angle approximation $\sin(x)=x$.

 





\textbf{Solution 3 (Alternate Finish)}

Let S be the sum of the sequence. We begin the same as in Solution 1 to get$S\sin(1)=\cot(45)-\cot(46)+\cot(47)-\cot(48)+...+\cot(133)-\cot(134)$. Observe that this "almost telescopes," if only we had some extra terms. Consider adding the sequence $\frac{1}{\sin(46)\sin(47)}+\frac{1}{\sin(48)\sin(49)}+...+\frac{1}{\sin(134)\sin(135)}$. By the identity $\sin(x)=\sin(180-x)$, this sequence is equal to the original one, simply written backwards. By the same logic as before, we may rewrite this sequence as$S\sin(1)=\cot(46)-\cot(47)+\cot(48)-\cot(49)+...+\cot(134)-\cot(135)$,and when we add the two sequences, they telescope to give $2S\sin(1)=\cot(45)-\cot(135)=2$.Hence, $S=\frac{1}{\sin(1^\circ)}$, and our angle is $\boxed{001}$.

 





\textbf{Solution 4}

First, multiply $\sin n^{\circ}$ on both sides.
 \begin{align*}     \frac{\sin n^\circ}{\sin m^\circ \sin (m+1)^\circ}     &= \frac{\sin (k+n-k)^\circ}{\sin m^\circ \sin (m+1)^\circ} \\     &= \frac{\sin (k+n)^\circ \cos k^\circ - \sin k^\circ \cos (k+n)^\circ}{\sin m^\circ \sin (m+1)^\circ} \end{align*}
Let $k = m$ since $k$ is could be any number.
 \begin{align*}     &\quad \ \frac{\sin (k+n)^\circ \cos k^\circ}{\sin m^\circ \sin (m+1)^\circ} - \frac{\sin k^\circ \cos (k+n)^\circ}{\sin m^\circ \sin (m+1)^\circ} \\[0.5em]     &=\frac{\sin (m+n)^\circ \cos m^\circ}{\sin m^\circ \sin (m+1)^\circ} - \frac{\sin m^\circ \cos (m+n)^\circ}{\sin m^\circ \sin (m+1)^\circ} \end{align*}


 \textbf{Lemma} \begin{verbatim}
$n$ is equal to $1$.
\end{verbatim}\textbf{Proof} 
 \begin{align*}     &=\frac{\sin (m+1)^\circ \cos m^\circ}{\sin m^\circ \sin (m+1)^\circ} - \frac{\sin m^\circ \cos (m+1)^\circ}{\sin m^\circ \sin (m+1)^\circ} \\[0.5em]     &=\frac{\cos m^\circ}{\sin m^\circ} - \frac{\cos (m+1)^\circ}{\sin (m+1)^\circ} \\[0.5em]     &= \cot m^\circ - \cot (m+1)^\circ \end{align*}
The sum of all numbers could be written. Moreover, notice that $\cot\alpha + \cot\beta = 0$ if $\alpha + \beta = 180^\circ$.
 \begin{align*}     &\quad \ \cot 45^\circ - \cot 46^\circ + \cot 47^\circ - \cot 48^\circ + \cdots - \cot 132^\circ + \cot 133^\circ - \cot 134^\circ \\     &= (\cot 45^\circ + \cot 47^\circ + \cdots + \cot 89^\circ + \cot 91^\circ + \dots + \cot 133^\circ) \\     &\qquad\qquad\qquad\qquad\qquad\qquad - (\cot 46^\circ + \dots + \cot 88^\circ + \cot 90^\circ + \cot 92^\circ + \dots + \cot 134^\circ) \\     &= \cot 45^\circ - \cot 90^\circ \\     &= 1 \end{align*}
Because $1 = 1$, the lemma is true.

 Q.E.D.

 $n$ could be 1. Moreover, there are no smaller positive integer less than 1 to test. Thus, the least positive integer $n$ that satisfies the given condition is $\boxed{001}$.

 


 




\end{enumerate}